\section{Wisconsin, Georgia, and Pennsylvania Results}
\label{sec:wi-ga-pa}

\subsection{Wisconsin ($k = 8$)}

Wisconsin's $k = 8 = 2^3$ seat count produces a three-level binary tree:
state $\to$ 2 halves $\to$ 4 quarters $\to$ 8 districts.  This hierarchical
structure tends to produce compact districts aligned with the state's
east-west geography.

\paragraph{Compactness.}
The \ar{} plan achieves $\bar\pp_{\mathrm{AR}}^{\mathrm{WI}} = 0.388$,
placing it at the $\mathbf{72\text{nd percentile}}$ of the DeFord ensemble.

\paragraph{Partisan outcome.}
The \ar{} plan produces 4D/4R districts, at the $\mathbf{61\text{st
percentile}}$ of the Democratic seat distribution.  The ensemble median
for Wisconsin is also 4D/4R under the 2020 presidential baseline, reflecting
the highly competitive nature of the state.

\paragraph{Minority-majority districts.}
At $k = 8$ with the state's demographic composition, the expected number of
minority-majority districts is approximately 1 (the Milwaukee-area district).
The \ar{} plan produces 1 MM district, at the $58$th percentile of the
ensemble MM distribution.

\begin{center}
\begin{tabular}{lccc}
\toprule
Statistic & AR value & Ensemble median & AR percentile \\
\midrule
Mean PP         & 0.388 & 0.301 & 72nd \\
Democratic seats & 4     & 4     & 61st \\
MM districts    & 1     & 1     & 58th \\
\bottomrule
\end{tabular}
\end{center}

\subsection{Georgia ($k = 14$)}

Georgia has the same $k = 14 = 2 \times 7$ structure as North Carolina but
a very different geographic partisan landscape.  Democratic voters are
concentrated in the Atlanta metropolitan area (Fulton, DeKalb, Gwinnett,
Clayton counties) and the Black Belt counties in the southwestern part of the
state.  Republican voters dominate a large contiguous rural region.

\paragraph{Compactness.}
The \ar{} plan achieves $\bar\pp_{\mathrm{AR}}^{\mathrm{GA}} = 0.395$,
placing it at the $\mathbf{65\text{th percentile}}$ of the DeFord ensemble.
Georgia is slightly below the NC and WI compactness percentiles because the
state's irregular shape (multiple peninsular geographic features) makes
compact partitioning harder.

\paragraph{Partisan outcome.}
The \ar{} plan produces 5D/9R districts, at the $\mathbf{38\text{th
percentile}}$ of the Democratic seat distribution.  This is the most
notable outlier in our study: the \ar{} plan is more Republican-leaning
than 62\% of ensemble plans.

This result requires explanation.  It does \emph{not} indicate algorithmic
bias toward Republicans.  Rather, it reflects the geographic structure of
Georgia: the minimum-edge-cut partition of the census-tract graph tends to
produce a small number of large, compact Democratic districts in the Atlanta
area and a larger number of smaller Republican districts covering the rural
and suburban areas.  This is the Rodden effect \citep{rodden2019why}: the
geographic concentration of Democratic voters in cities reduces the number
of districts Democrats can efficiently win under any compact plan.

The 38th percentile means that 38\% of \emph{random} valid plans also
produce 5 or fewer Democratic seats.  It is notable, but it is not an extreme
outlier in the way that a 99th or 1st percentile result would be.

\paragraph{Minority-majority districts.}
The \ar{} plan produces 5 MM districts in Georgia (reflecting the Black
Belt and Atlanta-area concentrations), at the $62$nd percentile of the
ensemble MM distribution.

\begin{center}
\begin{tabular}{lccc}
\toprule
Statistic & AR value & Ensemble median & AR percentile \\
\midrule
Mean PP         & 0.395 & 0.312 & 65th \\
Democratic seats & 5     & 6     & 38th \\
MM districts    & 5     & 5     & 62nd \\
\bottomrule
\end{tabular}
\end{center}

\subsection{Pennsylvania ($k = 17$)}

Pennsylvania's $k = 17$ is prime, which is unusual in the \ar{} framework:
there is no factorization tree, and the plan is computed as a single 17-way
minimum-cut partition of the state census-tract graph.  METIS handles
arbitrary-$k$ partitions, but the prime-$k$ case produces a different
geometric character than composite-$k$ plans.

\paragraph{Compactness.}
The \ar{} plan achieves $\bar\pp_{\mathrm{AR}}^{\mathrm{PA}} = 0.371$,
placing it at the $\mathbf{61\text{st percentile}}$.  This is the lowest
compactness percentile among the four directly-compared states.  The
direct 17-way partition is less compact than a hierarchical partition
because METIS must balance 17 regions simultaneously without the guidance
of the factorization tree.

\paragraph{Partisan outcome.}
The \ar{} plan produces 8D/9R districts, at the $\mathbf{52\text{nd
percentile}}$ of the Democratic seat distribution.  This is very close to
the median, reflecting Pennsylvania's status as a closely divided state.

\paragraph{Minority-majority districts.}
The \ar{} plan produces 3 MM districts (Philadelphia, Pittsburgh, and a
Lehigh Valley majority-Hispanic district), at the $59$th percentile of the
ensemble MM distribution.

\begin{center}
\begin{tabular}{lccc}
\toprule
Statistic & AR value & Ensemble median & AR percentile \\
\midrule
Mean PP         & 0.371 & 0.306 & 61st \\
Democratic seats & 8     & 8     & 52nd \\
MM districts    & 3     & 3     & 59th \\
\bottomrule
\end{tabular}
\end{center}
